\section{Introduction}
최근 구조화된 추론 프롬프팅 연구에서는, 문제 해결에 필요한 핵심적인 사고의 흐름을 간결하게 제시함으로써 모델의 추론 과정을 유도하는 방식이 반복적으로 관찰되어 왔다. 본 논문에서는 이러한 추론 흐름의 명시적 표현을 스켈레톤(skeleton)이라 통칭한다. 일반적으로 스켈레톤은 입력 질의에 대해 문제 풀이의 방향이나 전략을 고수준에서 선행 제시하고, 이를 바탕으로 이후 상세한 추론과 최종 답변이 생성되는 구조를 가진다.
% 언어모델의 스켈레톤(skeleton) 기반 추론은 문제 해결에 필요한 핵심적인 사고의 흐름을 간결하게 제시함으로써 모델의 추론 과정을 유도하는 방법으로 주목받고 있다. 일반적으로 스켈레톤은 입력 질의에 대해 문제 풀이의 방향이나 전략을 먼저 제시하고, 이후 상세한 추론과 최종 답변을 생성하는 방식으로 구성된다.

스켈레톤 기반 추론의 중요한 특징 중 하나는 추가 학습 없이(\emph{training-free}) 적용 가능하%다는 점이다. 다국어 환경에서 특정 언어에 대해 미세조정이나 추가 학습을 수행하는 것은 데이터 부족, 비용, 인프라 제약으로 인해 현실적으로 어려운 경우가 많다. 이러한 맥락에서 스켈레톤은 사전 학습된 모델의 추론 능력을 그대로 활용하면서도 추론 과정만을 조절할 수 있는 수단으로서, 
며 이는 저자원 언어 환경에서 실용적인 가능성을 가진다.

그러나 다국어 환경에서 스켈레톤을 실제로 적용할 때는 하나의 근본적인 설계 문제가 존재한다. 스켈레톤 기반 추론 파이프라인은 암묵적으로 세 가지 언어 선택을 포함한다. 즉, (1) 입력 질의의 언어, (2) 스켈레톤이 표현되는 언어, (3) 추론 및 답변이 생성되는 언어이다. 이 세 요소의 조합에 따라 모델이 활용하는 언어적 prior와 추론 경로는 달라질 수 있음에도, 기존 연구에서는 이러한 언어 조합의 설계 공간을 체계적으로 다루지 않았다.

기존의 스켈레톤 관련 연구들은 주로 영어 환경을 전제로, 스켈레톤의 구조나 프롬프트 형식을 어떻게 설계해야 성능이 향상되는지에 초점을 맞추어 왔다. 일부 연구에서는 스켈레톤을 다국어 또는 다양한 도메인에 적용하기도 했으나, 대부분 영어 스켈레톤을 고정적으로 사용하거나 언어 선택을 부차적인 요소로 취급하였다. 결과적으로, 실제 다국어 배치 환경에서 중요한 질문인 “어떤 언어로 스켈레톤을 사용하는 것이 적절한가”에 대해서는 충분한 분석이 이루어지지 않았다.

한편 최근의 선행 연구들은 추론이 수행되는 언어에 따라 동일한 모델이라도 성능이 달라질 수 있음을 보고하였으며, 질의 언어를 영어로 변환한 뒤 추론을 수행하는 방식이 효과적일 수 있음을 보여주었다. 이러한 결과들은 스켈레톤 또한 단순한 형식적 도구가 아니라, 모델이 사전 학습 과정에서 형성한 언어별 추론 prior를 활성화하는 매개체로 작동할 수 있음을 시사한다. 동시에, 스켈레톤의 효과가 언어 전반에 걸쳐 균일하게 나타나지 않을 가능성 또한 내포한다.

이러한 맥락에서 본 논문은 다음 질문을 제기한다. \textit{스켈레톤 기반 추론은 어떤 언어적 환경에서 효과적으로 작동하며, 입력 언어–스켈레톤 언어–추론 및 답변 언어의 조합은 추론 성능에 어떠한 영향을 미치는가? 특히 저자원 언어에서는 이러한 언어 조합의 영향이 어떻게 달라지는가?} 이 질문은 “어떤 스켈레톤이 좋은가”를 넘어서, ``스켈레톤을 어떤 언어로 사용하는 것이 바람직한가''라는 보다 근본적인 설계 문제에 해당한다.

이를 탐구하기 위해 본 논문에서는 \textbf{Language-Aware Skeleton Exploration}을 제안한다. 여기서 탐색은 새로운 알고리즘적 학습이나 파라미터 업데이트를 의미하지 않으며, 스켈레톤의 구조와 형식은 고정한 채 언어 조합에 따른 성능 변화를 체계적으로 분석하고, 효과적인 언어 구성을 식별하는 경험적 접근을 의미한다. 이러한 설정을 통해 우리는 스켈레톤 자체의 설계 요소와 언어 선택의 영향을 분리하여 분석할 수 있다.

% 여기서 `최적화'란 새로운 알고리즘적 학습이나 파라미터 업데이트를 의미하지 않으며, 스켈레톤의 구조와 형식은 고정한 채 언어 조합에 따른 성능 변화를 체계적으로 분석하고, 효과적인 언어 구성을 식별하는 경험적 접근을 의미한다. 이러한 설정을 통해 우리는 스켈레톤 자체의 설계 요소와 언어 선택의 영향을 분리하여 분석할 수 있다.

본 연구는 수학 추론이라는 상대적으로 스켈레톤 기반 접근이 안정적으로 작동하기 어렵다고 알려진 영역을 대상으로, 다양한 난이도의 벤치마크와 서로 다른 모델 규모, 그리고 고자원 언어부터 저자원 언어에 이르는 폭넓은 언어 설정에서 실험을 수행한다. 이를 통해 스켈레톤의 효과가 언어와 조건에 따라 어떻게 달라지는지를 체계적으로 분석하고, 다국어 추론 환경에서 스켈레톤을 보다 신뢰성 있게 활용하기 위한 실질적인 시사점을 도출한다.

요약하면, 본 논문의 기여는 다음과 같다.
\begin{itemize}
    \item 다국어 스켈레톤 기반 추론에서 입력 언어–스켈레톤 언어–추론 및 답변 언어의 조합을 핵심 설계 변수로 정식화한다.
    \item 추가 학습 없이 적용 가능한 언어 인식 스켈레톤 활용 프레임워크를 제안하고, 언어 조합에 따른 성능 변화를 체계적으로 분석한다.
    \item 수학 추론 벤치마크 전반에서 스켈레톤의 효과가 언어 및 조건에 따라 비균일하게 나타날 수 있음을 보이고, 저자원 언어 환경에서의 활용 가능성과 한계를 논의한다.
\end{itemize}
